% =====================================================================
% Parcial 3 - Análisis Espacial
% Distribución espacial del precio mayorista del tomate de árbol, 2024
% Compilar con pdflatex (dos veces). Las figuras van en la carpeta figuras/
% =====================================================================
\documentclass[aspectratio=169,11pt]{beamer}

% --- Idioma y codificación -------------------------------------------
\usepackage[spanish,es-noshorthands,es-tabla]{babel}
\usepackage[utf8]{inputenc}
\usepackage[T1]{fontenc}
\usepackage{lmodern}
\usepackage[scaled=0.92]{helvet}
\renewcommand{\familydefault}{\sfdefault}

% --- Paquetes ----------------------------------------------------------
\usepackage{booktabs}
\usepackage{multirow}
\usepackage{tikz}
\usetikzlibrary{shapes.geometric, positioning, calc, shadows}
\usepackage{fontawesome5}
\usepackage{amsmath}
\usepackage{array}
\usepackage{xcolor}
\usepackage{listings}
\usepackage{graphicx}

% --- Paleta: tomate de árbol, campo y mercado --------------------------
\definecolor{tomate}{HTML}{B5302A}      % rojo del tomate de árbol maduro
\definecolor{naranja}{HTML}{E07B39}     % naranja de la pulpa
\definecolor{hoja}{HTML}{3E7C3B}        % verde de la hoja
\definecolor{hojaclara}{HTML}{8DB36B}   % verde claro
\definecolor{tierra}{HTML}{6B4A2E}      % café de la tierra
\definecolor{crema}{HTML}{FFF8EE}       % fondo cálido
\definecolor{tinta}{HTML}{2E2A26}       % texto
\definecolor{gris}{HTML}{8A817A}

% --- Tema propio ------------------------------------------------------
\setbeamercolor{background canvas}{bg=crema}
\setbeamercolor{normal text}{fg=tinta}
\setbeamercolor{frametitle}{fg=tomate}
\setbeamercolor{structure}{fg=hoja}
\setbeamercolor{alerted text}{fg=tomate}
\setbeamercolor{block title}{bg=hoja, fg=white}
\setbeamercolor{block body}{bg=hoja!8}
\setbeamercolor{block title alerted}{bg=tomate, fg=white}
\setbeamercolor{block body alerted}{bg=tomate!8}
\setbeamercolor{block title example}{bg=naranja, fg=white}
\setbeamercolor{block body example}{bg=naranja!10}
\setbeamerfont{frametitle}{size=\Large, series=\bfseries}
\setbeamertemplate{navigation symbols}{}
\setbeamertemplate{blocks}[rounded][shadow=false]
\setbeamertemplate{itemize item}{\color{tomate}\small\faCircle}
\setbeamertemplate{itemize subitem}{\color{hoja}\scriptsize\faLeaf}
\setbeamertemplate{enumerate item}{\color{tomate}\bfseries\insertenumlabel.}

% Título de cada diapositiva con una franja verde-naranja debajo
\makeatletter
\setbeamertemplate{frametitle}{%
  \vspace{0.35cm}%
  \begin{beamercolorbox}[wd=\paperwidth, leftskip=0.6cm]{frametitle}%
    \usebeamerfont{frametitle}\insertframetitle%
    \ifx\insertframesubtitle\@empty\else\\[-1pt]{\normalsize\color{gris}\insertframesubtitle}\fi
  \end{beamercolorbox}%
  \vspace{-0.1cm}%
  \begin{tikzpicture}[remember picture, overlay]
    \fill[hoja] ($(current page.north west)+(0.6,-1.38)$) rectangle ++(2.2,0.07);
    \fill[naranja] ($(current page.north west)+(2.8,-1.38)$) rectangle ++(0.9,0.07);
  \end{tikzpicture}%
}
\makeatother

% Pie: sección a la izquierda, número a la derecha y barra de progreso
\setbeamertemplate{footline}{%
  \begin{tikzpicture}[remember picture, overlay]
    \fill[hoja!15] (current page.south west) rectangle ($(current page.south east)+(0,0.5)$);
    \fill[tomate] (current page.south west) rectangle
      ($(current page.south west)+({\paperwidth*(\insertframenumber/\inserttotalframenumber)},0.08)$);
    \node[anchor=west, text=hoja, font=\scriptsize] at ($(current page.south west)+(0.6,0.27)$)
      {\faSeedling\ \insertsectionhead};
    \node[anchor=east, text=gris, font=\scriptsize] at ($(current page.south east)+(-0.6,0.27)$)
      {\insertframenumber\,/\,\inserttotalframenumber};
  \end{tikzpicture}%
}

% Tomate de árbol dibujado (para portada y separadores)
\newcommand{\tomatearbol}[3]{% x, y, escala
  \begin{scope}[shift={($(current page.south west)+(#1,#2)$)}, scale=#3]
    \shade[ball color=tomate] (0,0) ellipse (0.55 and 0.75);
    \fill[white, opacity=0.35] (-0.2,0.3) ellipse (0.10 and 0.20);
    \draw[tierra, line width=1.2pt] (0,0.72) -- (0.05,1.05);
    \fill[hoja] (0.05,0.98) .. controls (0.5,1.25) and (0.75,1.0) .. (0.85,0.85)
                .. controls (0.5,0.80) and (0.25,0.85) .. (0.05,0.98);
  \end{scope}}

% Separador de sección
\AtBeginSection[]{%
  \begin{frame}[plain]
    \begin{tikzpicture}[remember picture, overlay]
      \fill[crema] (current page.south west) rectangle (current page.north east);
      \fill[hoja] (current page.south west) rectangle ($(current page.south west)+(0.45,\paperheight)$);
      \fill[naranja] ($(current page.south west)+(0.45,0)$) rectangle ($(current page.south west)+(0.6,\paperheight)$);
      \node[anchor=west, text=tomate, font=\Huge\bfseries] at ($(current page.west)+(1.3,0.3)$) {\insertsectionhead};
      \node[anchor=west, text=gris, font=\large] at ($(current page.west)+(1.3,-0.6)$)
        {Parte \insertsectionnumber};
      \tomatearbol{12.9}{4.6}{1.3}
      \tomatearbol{14.4}{3.7}{0.9}
    \end{tikzpicture}
  \end{frame}}

% Insertar figura si existe; si no, recuadro con el nombre del archivo.
% #3 = etiqueta: al hacer clic en la figura se salta a su versión ampliada.
\newlength{\altofig}
\newcommand{\figurabase}[2]{%
  \IfFileExists{figuras/#1}{\includegraphics[height=#2,keepaspectratio]{figuras/#1}}{%
    \setlength{\altofig}{#2}%
    \fcolorbox{gris}{white}{\parbox[c][\altofig][c]{1.17\altofig}{\centering\color{gris}\footnotesize
      \faImage\\[2pt]\texttt{figuras/\detokenize{#1}}}}}}
\newcommand{\figura}[3]{%
  \hyperlink{zoom-#3}{\figurabase{#1}{#2}}%
  \llap{\raisebox{3pt}{\colorbox{white}{\scriptsize\color{hoja}\faSearchPlus}}\hspace{3pt}}}

% Diapositiva de ampliación: figura a pantalla completa y botón para volver
\newcommand{\ampliacion}[3]{% archivo, etiqueta, título
  \begin{frame}[plain, label=zoom-#2]
    \vspace{-0.15cm}
    {\small\bfseries\color{tomate} #3}\hfill\hyperlink{origen-#2}{\beamerbutton{\faUndo\ Volver}}\par
    \vspace{0.1cm}
    \centering\figurabase{#1}{0.82\paperheight}
  \end{frame}}

% Código R
\lstset{language=R, basicstyle=\ttfamily\scriptsize, keywordstyle=\color{hoja}\bfseries,
  commentstyle=\color{gris}\itshape, stringstyle=\color{tomate}, backgroundcolor=\color{white},
  frame=single, rulecolor=\color{hoja!40}, breaklines=true, showstringspaces=false,
  literate={á}{{\'a}}1 {é}{{\'e}}1 {í}{{\'i}}1 {ó}{{\'o}}1 {ú}{{\'u}}1 {ñ}{{\~n}}1}


% Etiqueta "En R": indica en qué parte del script está el resultado
\newcommand{\enR}[1]{%
  \begin{tikzpicture}[remember picture, overlay]
    \node[anchor=north east, fill=hoja, text=white, rounded corners=3pt, font=\scriptsize\bfseries,
          inner sep=4pt] at ($(current page.north east)+(-0.5,-0.45)$) {\faCode\ #1};
  \end{tikzpicture}}

% Destacar cifras
\newcommand{\cifra}[1]{{\color{tomate}\bfseries #1}}
\newcommand{\sig}{\textcolor{hoja}{\faCheck}}
\newcommand{\nosig}{\textcolor{gris}{\faTimes}}

% --- Datos de portada -------------------------------------------------
\title{¿Más lejos de la cosecha, más caro?}
\subtitle{Distribución espacial del precio mayorista del tomate de árbol en Colombia, 2024}
\author{[Tu nombre]}
\institute{Análisis Espacial · Parcial 3 · Universidad Externado de Colombia}
\date{Semestre 2026-II}

\begin{document}

% =====================================================================
% PORTADA
% =====================================================================
\begin{frame}[plain]
  \begin{tikzpicture}[remember picture, overlay]
    \fill[crema] (current page.south west) rectangle (current page.north east);
    \fill[hoja] (current page.north west) rectangle ($(current page.north east)+(0,-0.25)$);
    \fill[naranja] ($(current page.north west)+(0,-0.25)$) rectangle ($(current page.north east)+(0,-0.35)$);
    \tomatearbol{12.6}{5.2}{1.5}
    \tomatearbol{14.4}{4.3}{1.1}
    \tomatearbol{13.4}{2.9}{0.8}
    \node[anchor=north west, text width=10.3cm, text=tomate, font=\huge\bfseries]
      at ($(current page.north west)+(0.9,-1.1)$) {\inserttitle};
    \node[anchor=north west, text width=10.3cm, text=tinta, font=\large]
      at ($(current page.north west)+(0.9,-3.3)$) {\insertsubtitle};
    \draw[naranja, line width=2pt] ($(current page.north west)+(0.9,-4.75)$) -- ++(3,0);
    \node[anchor=north west, text width=10cm, font=\normalsize, text=tinta]
      at ($(current page.north west)+(0.9,-5.0)$) {\textbf{\insertauthor}\\[2pt]
      {\small\color{gris}\insertinstitute\\ \insertdate}};
  \end{tikzpicture}
\end{frame}

% =====================================================================
\section{Pregunta y fundamento}
% =====================================================================
\begin{frame}{Pregunta, hipótesis y objetivos}
  \begin{alertblock}{\faQuestionCircle\ Pregunta}
    ¿El precio mayorista del tomate de árbol en 2024 presenta dependencia espacial entre
    municipios de Colombia, y se asocia con la cercanía a las zonas productoras?
  \end{alertblock}
  \vspace{0.15cm}
  \begin{columns}[T]
    \column{0.5\textwidth}
    \textbf{\color{hoja}Hipótesis}
    \begin{itemize}\small
      \item \textbf{H1:} municipios vecinos tienen precios parecidos (autocorrelación positiva).
      \item \textbf{H2:} más \textbf{distancia} a la gran producción $\Rightarrow$ \textbf{mayor} precio.
      \item \textbf{H3:} más \textbf{producción local} $\Rightarrow$ \textbf{menor} precio.
    \end{itemize}
    \column{0.5\textwidth}
    \textbf{\color{hoja}Objetivos}
    \begin{enumerate}\small
      \item Medir la autocorrelación global y local del precio.
      \item Estimar MCO y 5 modelos espaciales con 3 matrices $W$.
      \item Evaluar qué conclusiones son estables o sensibles.
    \end{enumerate}
  \end{columns}
\end{frame}

\begin{frame}{Fundamento teórico: el precio viaja con la cosecha}
  \begin{columns}[c]
    \column{0.58\textwidth}
    \begin{itemize}\small
      \item \textbf{Análisis espacial de precios} (Fackler y Goodwin, 2001): entre mercados
            conectados, los precios difieren a lo sumo en el \textbf{costo de transporte}.
      \item \textbf{Evidencia con SIPSA} (Banco de la República, 2020): los costos de
            transacción equivalen a cerca del \cifra{42\%} del precio en origen; \cifra{+10\%}
            en tiempo de viaje $\approx$ \cifra{+3\%} en ese costo.
      \item \textbf{Concentración productiva} (EVA, 2024): Cundinamarca y Antioquia producen
            el \cifra{70\%} del tomate de árbol del país.
    \end{itemize}
    \column{0.42\textwidth}
    \begin{tikzpicture}[node distance=0.55cm, every node/.style={font=\small, align=center}]
      \node[draw=hoja, fill=hoja!10, rounded corners, minimum width=4.2cm, minimum height=0.8cm] (p)
        {\faSeedling\ Zona productora};
      \node[draw=naranja, fill=naranja!10, rounded corners, minimum width=4.2cm, minimum height=0.8cm, below=of p] (t)
        {\faTruck\ Transporte e intermediación};
      \node[draw=tomate, fill=tomate!10, rounded corners, minimum width=4.2cm, minimum height=0.8cm, below=of t] (m)
        {\faStore\ Precio en el mercado};
      \draw[-latex, thick, tierra] (p) -- (t);
      \draw[-latex, thick, tierra] (t) -- (m);
    \end{tikzpicture}
  \end{columns}
  \vspace{0.2cm}
  \begin{exampleblock}{Relaciones esperadas}\small
    $X_1$ distancia a la producción: \textbf{positiva} \quad|\quad
    $X_2$ producción local: \textbf{negativa}
  \end{exampleblock}
\end{frame}

% =====================================================================
\section{Datos y matrices}
% =====================================================================
\begin{frame}{Base analítica: una fila = un municipio}
  \enR{Script 1}
  \begin{columns}[T]
    \column{0.52\textwidth}
    \textbf{\color{hoja}\faDatabase\ Fuentes}
    \begin{itemize}\small
      \item \textbf{SIPSA-P mensual 2024} (DANE): precios mayoristas, 54.042 registros,
            349 productos, 80 mercados.
      \item \textbf{EVA 2024} (MinAgricultura): producción de tomate de árbol en 269 municipios.
      \item \textbf{MGN 2025} y \textbf{DIVIPOLA} (DANE): polígonos y códigos.
    \end{itemize}
    \textbf{\color{hoja}\faMapMarker*\ Unidad de análisis}
    \begin{itemize}\small
      \item \cifra{41 municipios} con precio los 12 meses.
      \item Polígonos, MAGNA-SIRGAS Origen-Nacional (EPSG:9377), metros.
    \end{itemize}
    \column{0.48\textwidth}
    \begin{block}{\faCheckCircle\ Control de calidad}\small
      \begin{itemize}
        \item 0 faltantes y 0 duplicados producto–mercado–mes.
        \item Los archivos \texttt{.dta} y \texttt{.sav} traían el precio \textbf{mal
              convertido} (1.387 $\to$ 1,387): se usó solo el CSV.
        \item 80 mercados $\to$ 60 municipios; 3 nombres homologados con DIVIPOLA
              (Cali, Cúcuta, Cartagena).
        \item EVA: miles con coma (\texttt{10,935}); corregido para no perder 29 municipios.
      \end{itemize}
    \end{block}
  \end{columns}
\end{frame}

\begin{frame}{Cómo está organizado el código}
  \begin{columns}[T]
    \column{0.42\textwidth}
    \begin{block}{\faFileCode\ Script 1: base}\footnotesize
      \begin{itemize}
        \item Importar el CSV del SIPSA-P (UTF-8, \texttt{;}).
        \item Precio numérico y fechas.
        \item Mercado $\to$ municipio $\to$ código DIVIPOLA.
        \item Polígonos MGN y paso a EPSG:9377.
      \end{itemize}
      \textbf{Guarda}: base limpia, mapa de 60 municipios, fondo de Colombia.
    \end{block}
    \column{0.58\textwidth}
    \begin{block}{\faFileCode\ Script 2: análisis espacial}\footnotesize
      \begin{tabular}{@{}l l@{}}
        §2--3 & Cobertura de productos y variable $y$ \\
        §4    & Mapas iniciales \\
        §5    & Matrices W1, W2, W3 \\
        §7    & Autocorrelación global de $y$ \\
        §8    & LISA: Moran, Getis-Ord y Geary locales \\
        §9    & Variables $X_1$, $X_2$ y global de $X_1$ \\
        §10   & 16 modelos, diagnóstico, coeficientes, \\
              & impactos y mapas de residuos \\
      \end{tabular}
    \end{block}
  \end{columns}
  {\scriptsize\color{gris}\faCode\ La etiqueta verde de cada diapositiva indica la sección del script que produce ese resultado.}
\end{frame}

\begin{frame}{Variables del modelo}
  \enR{Script 2 · §3, §9, §10.1}
  \centering\small
  \begin{tabular}{>{\bfseries}l >{\raggedright\arraybackslash}p{5.6cm} l c}
    \toprule
    & Definición & Unidad & Esperado \\
    \midrule
    \color{tomate}$y$ & Mediana anual del precio mayorista del tomate de árbol & \$/kg & --- \\
    \color{hoja}$X_1$ & $\log(1+$ distancia al gran productor más cercano$)$ & km & $+$ \\
    \color{hoja}$X_2$ & $\log(1+$ producción local de tomate de árbol$)$ & t & $-$ \\
    \bottomrule
  \end{tabular}

  \vspace{0.35cm}
  \begin{columns}[T]
    \column{0.5\textwidth}
    \begin{itemize}\small
      \item \textbf{Mediana}: un mes atípico no distorsiona el precio típico del año.
      \item \textbf{Gran productor}: 10\% de municipios con más producción
            ($\geq$ 1.409 t), que aportan el \cifra{79\%} del total.
    \end{itemize}
    \column{0.5\textwidth}
    \begin{itemize}\small
      \item \textbf{Logaritmos}: el precio sube rápido en los primeros km y luego se aplana
            ($r$: 0,43 $\to$ \cifra{0,58}); 21 municipios no producen.
      \item \textbf{Multicolinealidad}: $r(X_1,X_2)=-0{,}65$, \textbf{VIF = 1,73}.
    \end{itemize}
  \end{columns}
\end{frame}

\begin{frame}{Mapas iniciales: dos clasificaciones de $y$}
  \hypertarget{origen-clasificaciones}{}
  \enR{Script 2 · §4}
  \begin{columns}[c]
    \column{0.62\textwidth}
    \centering\figura{mapa_y_clasificaciones.png}{5.6cm}{clasificaciones}
    \column{0.38\textwidth}
    \begin{itemize}\small
      \item Precio: de \cifra{\$1.615} (La Unión) a \cifra{\$4.027} (Florencia).
      \item \textbf{Patrón común}: bajo en Antioquia, alto en la periferia oriental y sur.
      \item \textbf{Cuantiles}: Pamplona, Ipiales, San Gil y Yopal parecen de los más caros.
      \item \textbf{Intervalos iguales}: solo 4 en el extremo (Florencia, Arauca, Cúcuta,
            Villavicencio).
    \end{itemize}
  \end{columns}
\end{frame}

\begin{frame}{Tres matrices de vecindad}
  \hypertarget{origen-w1}{}
  \enR{Script 2 · §5}
  \begin{columns}[c]
    \column{0.55\textwidth}\small
    \textbf{Queen no sirve}: los 41 municipios están dispersos y quedarían sin vecinos.
    \vspace{0.15cm}

    \begin{tabular}{l >{\raggedright\arraybackslash}p{4.6cm}}
      \toprule
      \color{tomate}\textbf{W1} & 4 vecinos más cercanos (\textbf{principal}).
                                  Todos con 4; red conectada. \\
      \color{hoja}\textbf{W2} & Umbral \cifra{247,8 km}: mínimo para que nadie quede aislado.
                                 De 1 (Arauca) a 21 vecinos; el Caribe queda separado. \\
      \color{naranja}\textbf{W3} & Vecinos de W2 con peso $1/d$: los cercanos pesan más. \\
      \bottomrule
    \end{tabular}

    \vspace{0.2cm}
    {\scriptsize Ej. Ibagué: Armenia (53 km) pesa 4,8\% en W2 y \cifra{11,8\%} en W3.
    Todas estandarizadas por filas; mismo orden que los datos.}
    \column{0.45\textwidth}
    \centering\figura{W1_vecinos.png}{5.8cm}{w1}
  \end{columns}
\end{frame}

\begin{frame}{W1 frente a W2: cómo se conectan los municipios}
  \enR{Script 2 · §5.3--5.4}
  \hypertarget{origen-redw1}{}\hypertarget{origen-redw2}{}
  \begin{columns}[T]
    \column{0.5\textwidth}
    \centering{\color{tomate}\bfseries W1: 4 vecinos más cercanos}\\[3pt]
    \figura{W1_red.png}{3.7cm}{redw1}
    \begin{itemize}\footnotesize
      \item Todos con 4 vecinos; \textbf{una sola red} conectada.
      \item La costa se une al interior por Antioquia.
    \end{itemize}
    \column{0.5\textwidth}
    \centering{\color{hoja}\bfseries W2: vecinos a menos de 247,8 km}\\[3pt]
    \figura{W2_red.png}{3.7cm}{redw2}
    \begin{itemize}\footnotesize
      \item \textbf{El Caribe queda separado} del interior.
      \item Centro andino denso: hasta \cifra{21} vecinos;
            Arauca con 1.
    \end{itemize}
  \end{columns}
\end{frame}

% =====================================================================
\section{Autocorrelación global}
% =====================================================================
\begin{frame}{Doce resultados globales}
  \enR{Script 2 · §7 y §9.3}
  \framesubtitle{Prueba $z$ bajo aleatorización, dos colas, $\alpha=0{,}05$}
  \centering\scriptsize
  \begin{tabular}{l l l r r r r c}
    \toprule
    Variable & Índice & $W$ & Estadístico & Esperado & $z$ & $p$ & Sig. \\
    \midrule
    \multirow{6}{*}{\color{tomate}\textbf{$y$: precio}}
      & I de Moran & W1 & 0,491 & $-0{,}025$ & 5,42 & $6{,}1\times10^{-8}$ & \sig \\
      & I de Moran & W2 & 0,229 & $-0{,}025$ & 3,50 & $4{,}7\times10^{-4}$ & \sig \\
      & C de Geary & W1 & 0,438 & 1 & 5,35 & $4{,}5\times10^{-8}$ & \sig \\
      & C de Geary & W2 & 0,692 & 1 & 3,94 & $4{,}0\times10^{-5}$ & \sig \\
      & G Getis-Ord & W1 & 0,0253 & 0,025 & 0,72 & 0,474 & \nosig \\
      & G Getis-Ord & W2 & 0,0250 & 0,025 & $-0{,}05$ & 0,964 & \nosig \\
    \midrule
    \multirow{6}{*}{\color{hoja}\textbf{$X_1$: distancia}}
      & I de Moran & W1 & 0,577 & $-0{,}025$ & 6,38 & $1{,}7\times10^{-10}$ & \sig \\
      & I de Moran & W2 & 0,474 & $-0{,}025$ & 6,93 & $4{,}1\times10^{-12}$ & \sig \\
      & C de Geary & W1 & 0,354 & 1 & 5,85 & $5{,}0\times10^{-9}$ & \sig \\
      & C de Geary & W2 & 0,516 & 1 & 5,96 & $2{,}5\times10^{-9}$ & \sig \\
      & G Getis-Ord & W1 & 0,0275 & 0,025 & 2,44 & 0,015 & \sig \\
      & G Getis-Ord & W2 & 0,0266 & 0,025 & 2,28 & 0,023 & \sig \\
    \bottomrule
  \end{tabular}
\end{frame}

\begin{frame}{Lectura global: los vecinos se parecen}
  \hypertarget{origen-dispersion}{}
  \enR{Script 2 · §7.1}
  \begin{columns}[c]
    \column{0.5\textwidth}
    \centering\figura{dispersion_moran_W1.png}{5.4cm}{dispersion}
    \column{0.5\textwidth}
    \begin{itemize}\small
      \item \textbf{Moran y Geary} coinciden: autocorrelación \textbf{positiva} con W1 y W2.
      \item \textbf{Más fuerte con W1}: I = 0,491 $\to$ 0,229 con W2. La similitud es
            \textbf{local}; la vecindad amplia la diluye.
      \item \textbf{G de $y$ no significativa}: coexisten grupos altos y bajos que se
            compensan.
      \item \textbf{$X_1$ sí concentra valores altos}: la periferia lejana a la producción.
    \end{itemize}
  \end{columns}
\end{frame}

% =====================================================================
\section{Patrones locales}
% =====================================================================
\begin{frame}{Moran local (W1): ¿dónde están los grupos?}
  \hypertarget{origen-moranlocal}{}
  \enR{Script 2 · §8.1--8.2}
  \begin{columns}[c]
    \column{0.48\textwidth}
    \centering\figura{mapa_moran_local.png}{5.8cm}{moranlocal}
    \column{0.52\textwidth}\small
    \begin{itemize}
      \item \cifra{8 Bajo-Bajo}: Medellín, La Ceja, La Unión, Peñol, Rionegro,
            San Vicente Ferrer, El Santuario, Sonsón.
      \item \cifra{1 Alto-Alto}: Arauca (\$3.889). \quad \cifra{1 Bajo-Alto}: Bucaramanga.
      \item 31 no significativos ($p\geq0{,}05$).
    \end{itemize}
    \begin{exampleblock}{Verificación: El Santuario}\footnotesize
      Precio \$1.753 ($z=-1{,}71$); vecinos (La Unión, Peñol, Rionegro, San Vicente)
      $\approx$ \$1.900 (rezago $-1{,}47$). $I_i = 2{,}58$, $p = 0{,}001$ $\Rightarrow$
      \textbf{Bajo-Bajo significativo}.
    \end{exampleblock}
  \end{columns}
\end{frame}

\begin{frame}{Getis-Ord y Geary locales}
  \hypertarget{origen-getis}{}\hypertarget{origen-geary}{}
  \enR{Script 2 · §8.3--8.4}
  \begin{columns}[T]
    \column{0.5\textwidth}
    \centering\figura{mapa_getis_local.png}{3.5cm}{getis}\\[2pt]
    {\small \cifra{8 puntos fríos} (Antioquia) y \cifra{2 calientes} (oriente).}
    \column{0.5\textwidth}
    \centering\figura{mapa_geary_local.png}{3.5cm}{geary}\\[2pt]
    {\small Mayor contraste: \cifra{Túquerres} (4,35) y \cifra{Florencia} (4,19).}
  \end{columns}
  \vspace{0.2cm}
  \begin{alertblock}{Global vs. local}\small
    El patrón global nace de un grupo de precios bajos en la \textbf{zona productora} antioqueña.
    41 pruebas a la vez: $\approx 2$ falsos positivos esperados al 5\%; se hallaron 10.
  \end{alertblock}
\end{frame}

% =====================================================================
\section{Modelos y diagnóstico}
% =====================================================================
\begin{frame}{Seis especificaciones, 16 ajustes}
  \enR{Script 2 · §10.2--10.3}
  \centering\small
  \begin{tabular}{>{\bfseries\color{hoja}}l l l}
    \toprule
    Modelo & Ecuación & Componente espacial \\
    \midrule
    MCO   & $y = X\beta + \varepsilon$ & Ninguno (se estima una vez) \\
    SLM   & $y = \rho Wy + X\beta + \varepsilon$ & Precio de los vecinos \\
    SEM   & $y = X\beta + u,\ u = \lambda Wu + \varepsilon$ & Error espacial \\
    SLX   & $y = X\beta + WX\theta + \varepsilon$ & $X$ de los vecinos \\
    SDM   & $y = \rho Wy + X\beta + WX\theta + \varepsilon$ & Precio y $X$ de los vecinos \\
    SARAR & $y = \rho Wy + X\beta + u,\ u = \lambda Wu + \varepsilon$ & Precio y error \\
    \bottomrule
  \end{tabular}

  \vspace{0.3cm}
  \begin{block}{MCO de referencia}\small
    $\hat\beta_{X_1} =$ \cifra{150,2} ($p=0{,}017$) \quad $\hat\beta_{X_2} = -46{,}0$ ($p=0{,}18$)
    \quad $R^2 = 0{,}365$ \quad Moran de residuos con W1: \cifra{0,158} ($p=0{,}025$)
    $\Rightarrow$ se justifican modelos espaciales.
  \end{block}
\end{frame}

\begin{frame}{Diagnóstico: 18 filas modelo–matriz}
  \enR{Script 2 · §10.4}
  \centering\scriptsize
  \begin{tabular}{l rrr rrr rrr}
    \toprule
    & \multicolumn{3}{c}{\color{tomate}\textbf{W1}} & \multicolumn{3}{c}{\color{hoja}\textbf{W2}}
    & \multicolumn{3}{c}{\color{naranja}\textbf{W3}} \\
    \cmidrule(lr){2-4}\cmidrule(lr){5-7}\cmidrule(lr){8-10}
    Modelo & pseudo-$R^2$ & $p$ Moran & AIC & pseudo-$R^2$ & $p$ Moran & AIC & pseudo-$R^2$ & $p$ Moran & AIC \\
    \midrule
    MCO   & 0,365 & \textbf{0,025} & 631 & 0,365 & \textbf{0,042} & \textbf{631} & 0,365 & \textbf{0,015} & 631 \\
    SLM   & 0,490 & 0,617 & \cifra{626} & 0,400 & 0,531 & 632 & 0,457 & 0,915 & \cifra{629} \\
    SEM   & 0,443 & 0,870 & 630 & 0,399 & 0,626 & 632 & 0,429 & 0,698 & 630 \\
    SLX   & 0,487 & 0,053 & 627 & 0,397 & \textbf{0,010} & 633 & 0,402 & \textbf{0,003} & 633 \\
    SDM   & 0,514 & 0,862 & 627 & 0,419 & 0,520 & 634 & 0,459 & 0,649 & 632 \\
    SARAR & 0,571 & 0,991 & \cifra{626} & 0,399 & 0,698 & 634 & 0,468 & 0,779 & 631 \\
    \bottomrule
  \end{tabular}

  \vspace{0.25cm}
  \begin{itemize}\small
    \item Varianza residual: de 240.465 (MCO) a 162.601 (SARAR-W1). En negrita: Moran residual significativo.
    \item SLM, SEM, SDM y SARAR \textbf{eliminan} la autocorrelación residual; SLX no.
    \item Menor AIC: \textbf{SLM/SARAR con W1}, \textbf{SLM con W3}; con \textbf{W2 gana el MCO}.
  \end{itemize}
\end{frame}

% =====================================================================
\section{Robustez}
% =====================================================================
\begin{frame}{Coeficientes de las dos $X$ centrales}
  \enR{Script 2 · §10.5}
  \centering\scriptsize
  \begin{tabular}{l r r r r r r}
    \toprule
    Ajuste & $X_1$ & $X_2$ & $\rho$ & $\lambda$ & $WX_1$ & $WX_2$ \\
    \midrule
    MCO      & \textbf{150,2} (0,017) & $-46{,}0$ (0,18) & & & & \\
    SLM-W1   & 84,0 (0,13) & $-36{,}4$ (0,20) & \textbf{0,47} (0,001) & & & \\
    SEM-W1   & 103,5 (0,12) & $-27{,}9$ (0,35) & & \textbf{0,46} (0,003) & & \\
    SLX-W1   & 11,7 (0,88) & $-55{,}6$ (0,09) & & & 138,2 (0,28) & $-72{,}6$ (0,29) \\
    SDM-W1   & 11,6 (0,87) & $-47{,}6$ (0,10) & 0,27 (0,15) & & 107,5 (0,36) & $-44{,}6$ (0,49) \\
    SARAR-W1 & 60,1 (0,18) & $-37{,}8$ (0,14) & \textbf{0,66} ($<$0,001) & $-0{,}62$ (0,14) & & \\
    \midrule
    SLM-W2   & 107,7 (0,08) & $-52{,}3$ (0,10) & 0,34 (0,11) & & & \\
    SEM-W2   & \textbf{128,8} (0,043) & $-52{,}6$ (0,08) & & 0,35 (0,14) & & \\
    SLX-W2   & 114,3 (0,19) & $-43{,}4$ (0,22) & & & 163,3 (0,23) & 94,4 (0,23) \\
    SDM-W2   & 101,4 (0,20) & $-46{,}7$ (0,14) & 0,29 (0,23) & & 112,9 (0,38) & 88,1 (0,22) \\
    SARAR-W2 & 114,2 (0,15) & $-53{,}7$ (0,09) & 0,22 (0,67) & 0,18 (0,77) & & \\
    \midrule
    SLM-W3   & 83,3 (0,15) & $-43{,}6$ (0,14) & \textbf{0,46} (0,006) & & & \\
    SEM-W3   & 115,0 (0,08) & $-41{,}0$ (0,17) & & \textbf{0,43} (0,022) & & \\
    SLX-W3   & 47,0 (0,61) & $-52{,}7$ (0,12) & & & 149,2 (0,23) & $-12{,}5$ (0,85) \\
    SDM-W3   & 37,4 (0,65) & $-47{,}5$ (0,11) & \textbf{0,40} (0,041) & & 84,1 (0,46) & $-2{,}0$ (0,97) \\
    SARAR-W3 & 75,4 (0,25) & $-43{,}9$ (0,13) & 0,53 (0,07) & $-0{,}14$ (0,80) & & \\
    \bottomrule
  \end{tabular}\\[2pt]
  {\tiny Coeficiente (valor $p$). En negrita: significativo al 5\%. $X$ en logaritmo.}
\end{frame}

\begin{frame}{Dos comparaciones de robustez}
  \enR{Script 2 · §10.5}
  \begin{columns}[T]
    \column{0.5\textwidth}
    \begin{block}{\faSync\ Cambiar $W$, misma familia (SLM)}\small
      \begin{itemize}
        \item $X_1$: 84 $\to$ 108 $\to$ 83. \textbf{Signo y magnitud estables}, no significativa.
        \item $\rho$: \cifra{0,47*} (W1) $\to$ 0,34 (W2) $\to$ \cifra{0,46*} (W3).
        \item La dependencia aparece con vecindades \textbf{locales} y se pierde con la amplia.
      \end{itemize}
    \end{block}
    \column{0.5\textwidth}
    \begin{block}{\faLayerGroup\ Cambiar familia, misma $W$ (W1)}\small
      \begin{itemize}
        \item MCO $\to$ SLM/SEM: $X_1$ baja de 150 a 84--103; $\rho$ y $\lambda$ captan la
              dependencia.
        \item SLX/SDM: $X_1 \approx 12$ y $WX_1 \approx 108$--138: la distancia propia y la
              de los vecinos casi coinciden (Moran de $X_1$ = 0,58).
      \end{itemize}
    \end{block}
  \end{columns}
\end{frame}

\begin{frame}{Residuos con los mismos cortes}
  \hypertarget{origen-residuos}{}
  \enR{Script 2 · §10.7}
  \begin{columns}[c]
    \column{0.55\textwidth}
    \centering\figura{mapa_residuos.png}{6.1cm}{residuos}
    \column{0.45\textwidth}\small
    \begin{itemize}
      \item SLM-W1 corrige errores agrupados del MCO (más municipios en blanco).
      \item SLM-W2 casi igual al MCO.
      \item Errores persistentes:
        \begin{itemize}
          \item \textbf{Subestimados}: Florencia (+959), Santa Marta (+826), Villavicencio (+735).
          \item \textbf{Sobreestimados}: Túquerres ($-944$), Manizales ($-732$), Popayán ($-676$).
        \end{itemize}
      \item Son los mismos municipios de mayor contraste en Geary local.
    \end{itemize}
  \end{columns}
\end{frame}

% =====================================================================
\section{Interpretación y conclusión}
% =====================================================================
\begin{frame}{Modelo elegido: SLM con W1}
  \enR{Script 2 · §10.6}
  \begin{columns}[T]
    \column{0.47\textwidth}
    \begin{block}{\faTrophy\ Por qué}\small
      \begin{itemize}
        \item \textbf{Vecindad} local y pareja, coherente con costos de transporte.
        \item \textbf{AIC} = 626 (MCO: 631).
        \item \textbf{Residuos} sin autocorrelación ($p = 0{,}617$).
        \item $\rho = 0{,}47$ significativo; el SARAR empata en AIC pero su $\lambda$ no es
              significativo.
      \end{itemize}
    \end{block}
    \column{0.53\textwidth}
    \textbf{\color{hoja}Impactos de $X_1$} (al duplicarse la distancia, $\times\ln 2$)\\[3pt]
    \centering\small
    \begin{tabular}{l r r}
      \toprule
      Efecto & Coef. & \$/kg \\
      \midrule
      Directo   & 88,8  & $\approx$ +62 \\
      Indirecto & 68,4  & $\approx$ +47 \\
      Total     & 157,2 & $\approx$ \cifra{+109} \\
      \bottomrule
    \end{tabular}\\[3pt]
    {\footnotesize $p$ entre 0,11 y 0,25: no significativos al 5\%. $X_2$: total $-68{,}2$ ($p=0{,}28$).}
  \end{columns}
  \vspace{0.15cm}
  \begin{exampleblock}{Qué añade frente al MCO}\small
    El efecto total de la distancia (157) es similar al MCO (150), pero el SLM muestra que
    \textbf{parte se transmite entre vecinos} y corrige la dependencia en los errores.
  \end{exampleblock}
\end{frame}

\begin{frame}{Dos conclusiones sobre el fenómeno}
  \centering\small
  \begin{tabular}{>{\raggedright\arraybackslash}p{3.9cm} >{\raggedright\arraybackslash}p{5.3cm} >{\raggedright\arraybackslash}p{3.4cm}}
    \toprule
    Conclusión & Evidencia & Clasificación \\
    \midrule
    \textbf{1.} Más lejos de la producción, más caro el tomate de árbol
      & $X_1$ positiva en los \textbf{16 ajustes} (11,6 a 150,2); significativa solo en MCO y SEM-W2
      & \textcolor{naranja}{\textbf{Signo estable, significancia sensible}} \\[4pt]
    \textbf{2.} El precio de un municipio depende del de sus vecinos
      & Moran 0,49 (W1); $\rho$/$\lambda$ significativos con W1 y W3, no con W2
      & \textcolor{naranja}{\textbf{Sensible a la matriz}} \\
    \bottomrule
  \end{tabular}

  \vspace{0.35cm}
  \begin{alertblock}{\faReply\ Respuesta a la pregunta}\small
    Sí hay dependencia espacial: un grupo de precios bajos en la zona productora antioqueña y
    precios altos en la periferia. La distancia a la producción tiene el signo esperado, pero
    una vez modelada la dependencia entre vecinos su efecto deja de ser significativo.
  \end{alertblock}
\end{frame}

\begin{frame}{Limitaciones}
  \begin{columns}[T]
    \column{0.5\textwidth}
    \begin{block}{\faRoute\ Distancia en línea recta}\small
      No refleja vías, relieve ni cierres. \textbf{Consecuencia}: Florencia y Villavicencio
      quedan subestimados por \$735--959/kg.
    \end{block}
    \column{0.5\textwidth}
    \begin{block}{\faUsers\ Tamaño de la muestra}\small
      41 municipios con datos todo el año. \textbf{Consecuencia}: efectos imprecisos; los
      impactos tienen el signo esperado pero no son significativos.
    \end{block}
  \end{columns}
  \vspace{0.2cm}
  \begin{block}{\faSeedling\ Solo grandes productores en $X_1$}\small
    Se ignora el abastecimiento desde productores medianos cercanos (p. ej. Túquerres,
    Manizales, Popayán son más baratos de lo predicho). Además, SIPSA cubre centrales
    mayoristas, no precios al consumidor.
  \end{block}
\end{frame}

\begin{frame}{Referencias}
  \scriptsize
  \begin{itemize}
    \item[] Anselin, L. (1995). Local indicators of spatial association—LISA.
            \textit{Geographical Analysis, 27}(2), 93--115.
    \item[] Banco de la República. (2020). \textit{Abastecimiento y costos de transacción en los
            mercados de alimentos de las principales ciudades de Colombia} (Borradores de
            Economía, n.º 1128). [Completar autores]
    \item[] Bivand, R., \& Piras, G. (2015). Comparing implementations of estimation methods
            for spatial econometrics. \textit{Journal of Statistical Software, 63}(18), 1--36.
    \item[] DANE. (2025). \textit{Sistema de Información de Precios y Abastecimiento del Sector
            Agropecuario (SIPSA-P), base mensual 2024} [Conjunto de datos].
    \item[] DANE. (2025). \textit{Marco Geoestadístico Nacional (MGN) 2025, nivel municipio}
            [Conjunto de datos].
    \item[] Fackler, P. L., \& Goodwin, B. K. (2001). Spatial price analysis. En B. L. Gardner
            \& G. C. Rausser (Eds.), \textit{Handbook of Agricultural Economics} (Vol. 1,
            pp. 971--1024). Elsevier.
    \item[] Geary, R. C. (1954). The contiguity ratio and statistical mapping.
            \textit{The Incorporated Statistician, 5}(3), 115--145.
    \item[] Getis, A., \& Ord, J. K. (1992). The analysis of spatial association by use of
            distance statistics. \textit{Geographical Analysis, 24}(3), 189--206.
    \item[] Ministerio de Agricultura y Desarrollo Rural. (2025). \textit{Evaluaciones
            Agropecuarias Municipales (EVA) 2019--2025, base agrícola} [Conjunto de datos].
            datos.gov.co.
    \item[] Moran, P. A. P. (1950). Notes on continuous stochastic phenomena.
            \textit{Biometrika, 37}(1/2), 17--23.
  \end{itemize}
\end{frame}

% =====================================================================
% RESPALDO
% =====================================================================
\appendix
\section*{Respaldo}

\begin{frame}[fragile]{Respaldo: bloque de R que construye W2}
  \enR{Script 2 · §5.4}
  \begin{columns}[T]
    \column{0.55\textwidth}
\begin{lstlisting}
# Distancia de cada municipio a su vecino
# más cercano (en metros)
nb_k1 <- knn2nb(knearneigh(coo, k = 1))
dist1 <- nbdists(nb_k1, coo)

# Umbral: la mayor de esas distancias
max_dist <- max(unlist(dist1))
max_dist / 1000          # 247.8 km

# Vecinos a menos del umbral
nb2 <- dnearneigh(coo, d1 = 0, d2 = max_dist)
W2  <- nb2listw(nb2, style = "W")
\end{lstlisting}
    \column{0.45\textwidth}\small
    \textbf{Entrada}: centroides de los 41 municipios en metros (EPSG:9377).\\[4pt]
    \textbf{Operación}: busca el vecino más cercano de cada uno y toma la
    \textbf{máxima} de esas distancias: el umbral mínimo para que nadie quede aislado.\\[4pt]
    \textbf{Salida}: lista de vecinos (1 a 21, mediana 12) y matriz estandarizada por filas.
  \end{columns}
\end{frame}

\begin{frame}{Respaldo: producción de tomate de árbol (EVA 2024)}
  \begin{columns}[T]
    \column{0.5\textwidth}\small
    \begin{tabular}{l r}
      \toprule
      Departamento & \% producción \\
      \midrule
      Cundinamarca & 36,8 \\
      Antioquia & 33,1 \\
      Tolima & 8,5 \\
      Boyacá & 8,2 \\
      Huila & 2,7 \\
      \bottomrule
    \end{tabular}
    \column{0.5\textwidth}\small
    \begin{itemize}
      \item 269 municipios productores, 203.720 t.
      \item Mayores: San Bernardo, Santa Rosa de Osos, Pandi, Cabrera, Donmatías.
      \item 20 de los 41 municipios del estudio producen; los 5 grandes
            (Sonsón, San Vicente Ferrer, Ibagué, Peñol, La Unión) coinciden con el
            grupo Bajo-Bajo.
    \end{itemize}
  \end{columns}
\end{frame}

% =====================================================================
% AMPLIACIONES (se llega con clic sobre cada figura)
% =====================================================================
\ampliacion{mapa_y_clasificaciones.png}{clasificaciones}{Mapas iniciales: cuantiles e intervalos iguales}
\ampliacion{W1_vecinos.png}{w1}{W1: 4 vecinos más cercanos}
\ampliacion{W1_red.png}{redw1}{W1: red de 4 vecinos más cercanos}
\ampliacion{W2_red.png}{redw2}{W2: red de vecinos a menos de 247,8 km}
\ampliacion{dispersion_moran_W1.png}{dispersion}{Diagrama de dispersión de Moran (W1)}
\ampliacion{mapa_moran_local.png}{moranlocal}{Moran local (W1)}
\ampliacion{mapa_getis_local.png}{getis}{Getis-Ord local (W1)}
\ampliacion{mapa_geary_local.png}{geary}{Geary local (W1)}
\ampliacion{mapa_residuos.png}{residuos}{Residuos: MCO y SLM con W1, W2 y W3}

\end{document}
